\documentclass[11pt]{article}

\usepackage[margin=1in]{geometry}
\usepackage{amsmath,amssymb,amsthm}
\usepackage{hyperref}
\usepackage{enumitem}

\hypersetup{
  colorlinks=true,
  linkcolor=blue,
  urlcolor=blue,
  citecolor=blue
}
\emergencystretch=2em

\newcommand{\ellone}{\ell_1}
\newcommand{\Lone}{\mathcal L_1}

\title{Literature-Implied Obstructions for the Complemented-Subspace Problem in \(L_1\)}
\author{Math discovery packet}
\date{June 21, 2026}

\begin{document}
\maketitle

\section{Status}

This is a lightweight \emph{literature-implied partial answer}, not a full
solution of the source question.

The source paper is Dirk Werner, \emph{Nigel Kalton's work in isometrical
Banach space theory}, arXiv:1103.3153.  In Section 5, ``Operators on
\(L_1\)'', immediately after Theorem 5.2, Werner writes that it is still open
whether a complemented infinite-dimensional subspace of \(L_1\) must be
isomorphic to \(L_1\) or \(\ell_1\).

The full problem remains open in the checked literature.  For example,
arXiv:1811.06571 calls it the famous problem whether every infinite-dimensional
complemented subspace of \(L_1\) is isomorphic to \(\ell_1\) or \(L_1\), and
arXiv:2310.02196 describes the separable \(L_1\)-space CSP as a long-standing
open problem.

\section{Literature-Implied Subcases}

The main supporting paper is
\[
\text{D. de Hevia, G. Martinez-Cervantes, A. Salguero-Alarcon, and
P. Tradacete,}
\]
\[
\text{\emph{A negative solution to the complemented subspace problem in
Banach lattices}, arXiv:2310.02196.}
\]
Its Corollary 2.3 states that if \(X\) is a complemented subspace of an
\(L_1\)-space and \(X\) is isomorphic to a Banach lattice, then \(X\) is
isomorphic to an \(L_1\)-space.  The proof combines the Lindenstrauss--Rosenthal
fact that complemented subspaces of \(L_1\)-spaces are \(\Lone\)-spaces with
Proposition 2.1 of the same paper, which says that a Banach lattice that is an
\(\Lone\)-space is lattice isomorphic to an \(L_1\)-space.

Consequently, in the separable \(L_1[0,1]\) setting of Werner's question, the
Banach-lattice subcase reduces to the standard separable \(L_1\)-space
classification: \(L_1[0,1]\), \(\ellone\), or a finite-dimensional
\(\ell_1^n\).  Since the source question is infinite-dimensional, this gives:
\[
\boxed{\text{Any counterexample must not be isomorphic to a Banach lattice.}}
\]

There is a second classical obstruction.  Lindenstrauss and Pelczynski's
Theorem 6.1 on \(\Lone\)-spaces says that the only \(\Lone\)-space with an
unconditional basis is \(\ellone\).  This theorem is quoted, for example, in
arXiv:1912.08449 and applied in arXiv:1905.13188 to complemented subspaces of
\(L_1\).  Since a complemented subspace of \(L_1\) is an \(\Lone\)-space, one
gets:
\[
\boxed{\text{A complemented subspace of \(L_1\) with an unconditional basis is
isomorphic to \(\ellone\).}}
\]

Thus any infinite-dimensional counterexample to Werner's question must be a
non-lattice \(\Lone\)-space and must have no unconditional basis.

\section{Failed Route}

The counterexample in arXiv:2310.02196 to the Banach-lattice CSP does not
appear to produce a counterexample for \(L_1\) by duality.  The paper notes
that the dual of the Plebanek--Salguero space \(\mathrm{PS}\) is a
1-complemented subspace of \(C(K_{\mathcal B})^*\), hence is linearly isometric
to \(\ell_1(\Gamma)\) for some set \(\Gamma\).  Therefore the adjoint side
falls into the known atomic \(L_1\) alternative rather than producing a new
complemented subspace of \(L_1[0,1]\).

\section{Search Evidence and Scope}

\begin{itemize}[leftmargin=*]
\item Run indexes were searched for arXiv:1103.3153, the source title,
Kalton/L1/complemented-subspace keywords, and close variants.  No prior
durable packet for this exact target was found.
\item Local source checks included arXiv:1811.06571, arXiv:1905.13188,
arXiv:1912.08449, and arXiv:2310.02196.
\item External arXiv/web checks on 2026-06-21 found no later full resolution
of the complemented-subspace problem for \(L_1[0,1]\).
\end{itemize}

This packet should be read as a structural narrowing of the search space, not
as a solution of Werner's problem.  A full counterexample, if it exists, must
evade both the Banach-lattice and unconditional-basis mechanisms above.

\section{References}

\begin{itemize}[leftmargin=*]
\item D. Werner, \emph{Nigel Kalton's work in isometrical Banach space
theory}, arXiv:1103.3153.
\item D. de Hevia, G. Martinez-Cervantes, A. Salguero-Alarcon, and
P. Tradacete, \emph{A negative solution to the complemented subspace problem
in Banach lattices}, arXiv:2310.02196.
\item J. Lindenstrauss and A. Pelczynski, \emph{Absolutely summing operators
in \(\mathcal L_p\)-spaces and their applications}, Studia Math. 29 (1968),
275--326.
\item F. Albiac, J. L. Ansorena, and P. Wojtaszczyk, \emph{On certain
subspaces of \(\ell_p\) for \(0<p\le 1\) and their applications to conditional
quasi-greedy bases in \(p\)-Banach spaces}, arXiv:1912.08449.
\item M. Novotny, \emph{Some Remarks on Schauder Bases in Lipschitz Free
Spaces}, arXiv:1905.13188.
\item W. B. Johnson, G. Schechtman, and T. Schlumprecht, \emph{Ideals in
\(L(L_1)\)}, arXiv:1811.06571.
\end{itemize}

\end{document}
